% ============================================================
% 第三章：基于 CTMC 的离散扩散模型基本原理
% 编译方式: xelatex -> bibtex -> xelatex -> xelatex
% ============================================================
\documentclass[12pt,a4paper]{ctexbook}

% ---------- 导言区 ----------
\usepackage{amsmath,amssymb,amsthm}
\usepackage{graphicx}
\usepackage{float}
\usepackage{subcaption}
\usepackage{caption}
\usepackage{booktabs}
\usepackage{array}
\usepackage{placeins}
\usepackage{natbib}
\usepackage{hyperref}
\usepackage{bm}

% 公式编号：阿拉伯数字（章节.序号）
\renewcommand{\theequation}{\arabic{chapter}.\arabic{equation}}

% 公式编号右对齐，Times New Roman 10.5pt
\makeatletter
\renewcommand{\maketag@@@}[1]{\hbox{\m@th\normalfont\fontsize{10.5}{12.6}\selectfont\fontfamily{ptm}\selectfont (#1)\quad}}
\makeatother

% 图表编号
\renewcommand{\thefigure}{\thechapter.\arabic{figure}}
\renewcommand{\thetable}{\thechapter.\arabic{table}}

% ---------- 定理类环境 ----------
\newtheorem{theorem}{定理}[chapter]
\newtheorem{lemma}[theorem]{引理}
\newtheorem{proposition}[theorem]{命题}
\newtheorem{corollary}[theorem]{推论}

\theoremstyle{definition}
\newtheorem{definition}{定义}[chapter]
\newtheorem{example}{例题}[chapter]
\newtheorem{remark}{注解}[chapter]

\renewenvironment{proof}[1][证明]{\par\noindent{\bfseries #1. }}{\hfill$\square$\par}

% ---------- 正文 ----------
\begin{document}

\chapter{基于 CTMC 的离散扩散模型基本原理}

连续时间马尔可夫链（Continuous-Time Markov Chain, CTMC）为离散状态空间的扩散模型提供了一个自然的数学框架。
本章将从基本的 CTMC 理论出发，逐一建立前向噪声过程、反向生成过程、连续时间 ELBO、
高效采样算法，最终统一离散时间与连续时间两种框架。

% ============================================================
\section{CTMC 回顾}
\label{sec:ctmc-basics}

在建立连续时间离散扩散模型之前，有必要回顾连续时间马尔可夫链的基本理论。
这些内容既是随机过程课程的核心知识，也是后续所有推导的数学基础。
本节将从扩散模型的实际需求出发，系统梳理 CTMC 的定义、基本方程、转移概率的闭式表达，以及时间反转等关键概念。

\subsection{数学记号}

\textbf{状态空间与数据。}考虑一个有限离散状态空间 $\mathcal{S} = \{1, 2, \dots, K\}$，
其中 $K$ 为状态总数。一条具有 $D$ 个维度的数据样本表示为
$\boldsymbol{x} = (x^1, x^2, \dots, x^D) \in \mathcal{S}^D$。
时刻 $t$ 的随机变量记为 $\boldsymbol{x}_t$，其中 $t$ 为连续时间，取值范围为 $[0, T]$，
$T$ 为终止时刻（通常取 $T = 1$）。

\textbf{概率分布。}用 $q_t(\boldsymbol{x})$ 表示前向过程在时刻 $t$ 的边缘概率分布，
$q_{t|s}(\boldsymbol{x}_t \mid \boldsymbol{x}_s)$ 表示转移核。
$q_{t|0}(\boldsymbol{x}_t \mid \boldsymbol{x}_0)$ 表示从初始干净数据出发经过时间 $t$ 后的分布。
模型分布的记号采用 $p_\theta$，其中 $\theta$ 为可学习参数。
神经网络输出的去噪预测表示为 $p_{0|t}^\theta(\boldsymbol{x}_0 \mid \boldsymbol{x}_t)$，
这一参数化方式源自 \cite[§2]{austin2021d3pm}，在连续时间框架中沿用自 \cite[§3.1]{campbell2022ctmc}。

\textbf{矩阵基础。}速率矩阵（生成元矩阵）记为 $R_t \in \mathbb{R}^{K \times K}$，
$[R_t]_{ij} = R_t(i, j)$ 表示时刻 $t$ 从状态 $i$ 到 $j$ 的瞬时速率
\cite[§3.1]{campbell2022ctmc}。
离散时间单步转移矩阵记为 $Q_t$，$[Q_t]_{ij} = q(x_t = j \mid x_{t-1} = i)$
\cite[§2]{austin2021d3pm}。
累积转移矩阵 $\overline{Q}_{t} = Q_1 Q_2 \cdots Q_t$ 刻画多步转移。

\textbf{噪声调度。}连续时间噪声调度函数记为 $\beta(t)$，
累积积分 $\overline{\beta}_{t|s} = \int_s^t \beta(\tau) \, d\tau$
\cite[§2.1]{zhao2025ud3}。
离散时间噪声参数 $\beta_t$ 与连续时间的关联为
$\beta_t = 1 - \exp\bigl(-\int_{t-1}^t \beta(\tau)\,d\tau\bigr)$
\cite[§4.1]{campbell2022ctmc}。

\textbf{基本向量与矩阵。}$\mathbf{1}$ 表示全 1 列向量，
$\mathbf{e}_i$ 表示第 $i$ 个标准基向量，$\delta_{ij}$ 为 Kronecker delta 符号，
$\odot$ 表示逐元素乘积，$\langle \cdot, \cdot \rangle$ 表示向量内积。
分类分布记为 $\text{Cat}(\boldsymbol{x}; \boldsymbol{p})$。
所有转移矩阵使用行向量左乘约定：
$\boldsymbol{x}_t \sim \text{Cat}(\boldsymbol{x}_{t-1} Q_t)$。

\subsection{连续时间马尔可夫链的基本定义}

连续时间马尔可夫链 $(\boldsymbol{x}_t)_{t \ge 0}$ 是一个在有限状态空间 $\mathcal{S}^D$ 上取值的随机过程，满足马尔可夫性：
\begin{equation}
\mathbb{P}(\boldsymbol{x}_t = \boldsymbol{y} \mid \{\boldsymbol{x}_u : u \le s\}) = \mathbb{P}(\boldsymbol{x}_t = \boldsymbol{y} \mid \boldsymbol{x}_s). \label{eq:markov-property}
\end{equation}
与离散时间马尔可夫链不同，CTMC 的状态可以在任意连续时刻发生变化，
这一特性使其天然适合描述连续的噪声注入过程。

\textbf{速率矩阵的定义与性质。}CTMC 的局部动力学完全由速率矩阵 $R_t$ 刻画。
对于 $i \neq j$，$R_t(i, j) \ge 0$；对角线元素定义为
\begin{equation}
R_t(i, i) = -\sum_{j \neq i} R_t(i, j), \label{eq:rate-diagonal}
\end{equation}
以保证行和为零。速率矩阵在无限小时间间隔上的含义由 \cite[§3.1]{campbell2022ctmc} 给出：
\begin{equation}
q_{t+\Delta t \mid t}(j \mid i) = \delta_{ij} + R_t(i, j) \Delta t + o(\Delta t), \quad \Delta t \to 0. \label{eq:infinitesimal}
\end{equation}

\subsection{Kolmogorov 方程}

CTMC 的边缘分布服从 Kolmogorov 前向方程 \cite[Appendix A]{campbell2022ctmc}：
\begin{equation}
\frac{\partial}{\partial t} q_t(\boldsymbol{x}) = \sum_{\boldsymbol{y}} q_t(\boldsymbol{y}) R_t(\boldsymbol{y}, \boldsymbol{x}). \label{eq:kolmogorov-fwd}
\end{equation}
矩阵形式（$q_t$ 为行向量）为
\begin{equation}
\frac{d}{dt} q_t = q_t R_t. \label{eq:kolmogorov-fwd-mat}
\end{equation}
对应的 Kolmogorov 后向方程为
\begin{equation}
\frac{\partial}{\partial t} q_{t|s}(\boldsymbol{x}_t \mid \boldsymbol{x}_s) = -\sum_{\boldsymbol{y}} R_t(\boldsymbol{x}_t, \boldsymbol{y}) \, q_{t|s}(\boldsymbol{y} \mid \boldsymbol{x}_s). \label{eq:kolmogorov-bwd}
\end{equation}
前向方程描述噪声如何破坏数据结构，后向方程则在反向过程的推导中发挥关键作用。

\subsection{转移概率与矩阵指数}

对于时间齐次的 CTMC，转移概率矩阵可解析地表示为矩阵指数：
\begin{equation}
P(t) = \exp(tR) = \sum_{n=0}^{\infty} \frac{(tR)^n}{n!}. \label{eq:matrix-exp}
\end{equation}
扩散模型中常使用可分离结构 \cite[§4.1]{campbell2022ctmc}\cite[§2.1]{zhao2025ud3}：
\begin{equation}
R_t = \beta(t) R_b, \label{eq:separable-rate}
\end{equation}
其中 $\beta(t) > 0$ 为噪声调度函数，$R_b$ 为基速率矩阵。在此设定下转移概率有闭式表达：
\begin{equation}
q_{t|s}(\boldsymbol{x}_t \mid \boldsymbol{x}_s) = \exp\!\Bigl( \overline{\beta}_{t|s} R_b \Bigr), \quad
\overline{\beta}_{t|s} = \int_s^t \beta(\tau)\, d\tau. \label{eq:closed-form-trans}
\end{equation}

\subsection{平稳分布与可逆性}

对于不可约非周期 CTMC，当 $t \to \infty$ 时 $q_t$ 收敛到平稳分布 $\pi$，满足 $\pi R = 0$。
若 CTMC 满足细致平衡条件
\begin{equation}
\pi(i) R(i, j) = \pi(j) R(j, i), \quad \forall i \neq j, \label{eq:detailed-balance}
\end{equation}
则称该链关于 $\pi$ 可逆。在生成模型的时间反转构造中，可逆性扮演了基础性角色。

\subsection{与离散时间马尔可夫链的关系}

CTMC 与离散时间马尔可夫链之间的联系 \cite[Appendix A]{austin2021d3pm}：
\begin{equation}
Q = \exp(\Delta t R) = I + R \Delta t + o(\Delta t). \label{eq:embedding}
\end{equation}
一阶近似 $Q \approx I + R \Delta t$ 对应了连续过程的欧拉离散化。
在 D3PM 等模型中 $Q_t$ 被直接参数化，在连续时间框架中则直接定义 $R_t$。
两种视角的统一构成了 \ref{sec:ud3} 节的核心主题。

\subsection{时间反转}

CTMC 的时间反转过程仍是 CTMC。设正向过程 $(\boldsymbol{x}_t)_{t \in [0, T]}$，
反向过程 $\bar{\boldsymbol{x}}_t = \boldsymbol{x}_{T - t}$ 的速率矩阵由
\cite[§3.1 Proposition 1]{campbell2022ctmc} 给出：
\begin{equation}
\hat{R}_t(\boldsymbol{x}, \tilde{\boldsymbol{x}}) = R_{T-t}(\tilde{\boldsymbol{x}}, \boldsymbol{x}) \frac{q_{T-t}(\tilde{\boldsymbol{x}})}{q_{T-t}(\boldsymbol{x})}, \quad \boldsymbol{x} \neq \tilde{\boldsymbol{x}}. \label{eq:time-reversal}
\end{equation}
当正向过程将数据破坏为噪声时，这个比率引导反向过程从噪声向数据恢复。
然而式 \eqref{eq:time-reversal} 中的边缘概率 $q_t$ 通常是不可知的，
这正是引入神经网络进行近似的根本原因。

% ============================================================
\section{前向过程}
\label{sec:forward}

前向过程的目标是在 $[0, T]$ 上通过一个 CTMC 将数据分布 $q_0$ 逐渐转化为易于处理的参考分布 $q_T$。

\subsection{可分离速率矩阵}

前向速率矩阵取可分离形式 \cite[§4.1]{campbell2022ctmc}\cite[§2.1]{zhao2025ud3}：
\begin{equation}
R_t = \beta(t) R_b. \label{eq:separable-forward}
\end{equation}
代入 \eqref{eq:closed-form-trans} 得到从 $s$ 到 $t$ 的转移矩阵：
\begin{equation}
q_{t|s}(\boldsymbol{x}_t \mid \boldsymbol{x}_s) = \exp\!\Bigl( \overline{\beta}_{t|s} R_b \Bigr). \label{eq:trans-s-to-t}
\end{equation}
当 $s = 0$ 时，从干净数据出发的边际转移为
\begin{equation}
q_{t|0}(\boldsymbol{x}_t \mid \boldsymbol{x}_0) = \exp\!\Bigl( \overline{\beta}_{t} R_b \Bigr). \label{eq:trans-t-to-0}
\end{equation}

\subsection{基速率矩阵的设计}

基速率矩阵 $R_b$ 的设计决定了噪声的微观结构，对应了三类转移矩阵的连续时间极限。

\textbf{均匀速率。}\cite[§3]{austin2021d3pm}\cite[§4.1]{campbell2022ctmc}：
\begin{equation}
R_b = \mathbf{1} \mathbf{1}^\top - K I, \label{eq:uniform-rate}
\end{equation}
其矩阵指数为
\begin{equation}
\exp(\alpha R_b) = e^{-\alpha K} I + \frac{1 - e^{-\alpha K}}{K} \mathbf{1} \mathbf{1}^\top. \label{eq:uniform-exp}
\end{equation}

\textbf{吸收态速率。}\cite[§3]{austin2021d3pm}\cite[§4.1]{campbell2022ctmc}：
引入吸收态 $m \in \mathcal{S}$（[MASK] token），
\begin{equation}
R_b = \mathbf{1} \mathbf{e}_m^\top - I, \label{eq:absorbing-rate}
\end{equation}
\begin{equation}
\exp(\alpha R_b) = e^{-\alpha} I + (1 - e^{-\alpha}) \mathbf{1} \mathbf{e}_m^\top. \label{eq:absorbing-exp}
\end{equation}

\textbf{离散化高斯速率。}\cite[§3]{austin2021d3pm}：
对于有序离散数据，
\begin{equation}
R_b(i, j) \propto \exp\!\Bigl( -\frac{|i - j|^2}{2\sigma^2} \Bigr), \quad i \neq j, \label{eq:gaussian-rate}
\end{equation}
保留了状态空间中的拓扑结构。

\subsection{维度分解}
\label{sec:dim-decomp}

全空间的速率矩阵维度为 $K^D \times K^D$，通过维度分解 \cite[§4.2 Proposition 3]{campbell2022ctmc}：
\begin{equation}
R_t^{1:D}(\tilde{\boldsymbol{x}}^{1:D}, \boldsymbol{x}^{1:D}) = \sum_{d=1}^D R_t^d(\tilde{x}^d, x^d) \;\delta_{\boldsymbol{x}^{1:D \setminus d}, \tilde{\boldsymbol{x}}^{1:D \setminus d}}. \label{eq:dim-decomp-rate}
\end{equation}
一次跳转只能改变一个维度，非零元数从 $K^D$ 降至 $D(K-1)+1$。
各维度之间前向独立：
\begin{equation}
q_{t|0}(\boldsymbol{x}_t \mid \boldsymbol{x}_0) = \prod_{d=1}^D q_{t|0}^d(x_t^d \mid x_0^d). \label{eq:dim-independence}
\end{equation}
这一独立性在反向过程中将被打破——去噪网络需要利用所有维度的信息进行联合预测。

\subsection{与离散时间前向过程的关系}

连续与离散框架的连接 \cite[Appendix A]{austin2021d3pm}\cite[§2.1]{zhao2025ud3}：
\begin{equation}
Q_t = \exp\!\Bigl( \beta_t R_b \Bigr), \quad \beta_t = \int_{t-1}^t \beta(\tau)\, d\tau. \label{eq:disc-conn}
\end{equation}
代入均匀速率和吸收态速率分别得到
\begin{align}
Q_t &= e^{-\beta_t K} I + \frac{1 - e^{-\beta_t K}}{K} \mathbf{1} \mathbf{1}^\top, \label{eq:uniform-disc} \\
Q_t &= e^{-\beta_t} I + (1 - e^{-\beta_t}) \mathbf{1} \mathbf{e}_m^\top, \label{eq:absorbing-disc}
\end{align}
与 D3PM 的离散转移矩阵形式完全一致。这一对应关系揭示了两种框架的内在统一性。

\subsection{噪声调度的选择}

离散时间框架中常用的调度包括线性调度、余弦调度和 Sohl-Dickstein 调度
\cite[§1]{austin2021d3pm}\cite[§4.1]{campbell2022ctmc}。
连续时间框架中调度函数 $\beta(t)$ 的灵活性是其相较于固定步长离散模型的一个重要优势。

% ============================================================
\section{反向过程与神经网络预测训练}
\label{sec:reverse}

给定前向 CTMC 的速率矩阵 $R_t$，需要回答两个问题：
精确的反向过程具有怎样的形式？如何用神经网络近似？

\subsection{精确反向速率}

由时间反转公式 \eqref{eq:time-reversal} 和 \cite[§3.1 Proposition 1]{campbell2022ctmc}，
在正向时间 $t$ 处的反向速率为
\begin{equation}
\hat{R}_t(\boldsymbol{x}, \tilde{\boldsymbol{x}}) = R_t(\tilde{\boldsymbol{x}}, \boldsymbol{x}) \frac{q_t(\tilde{\boldsymbol{x}})}{q_t(\boldsymbol{x})}, \quad \boldsymbol{x} \neq \tilde{\boldsymbol{x}}. \label{eq:reverse-exact}
\end{equation}
这一公式的推导可从 Kolmogorov 后向方程出发，详见 \cite[Appendix B.1]{campbell2022ctmc}。
然而边缘分布 $q_t$ 通常是未知且不可计算的，必须引入近似。

\subsection{基于去噪网络的反向速率参数化}

为获得可计算的反向速率，将 \eqref{eq:reverse-exact} 中的边缘概率比展开为
\begin{equation}
\frac{q_t(\tilde{\boldsymbol{x}})}{q_t(\boldsymbol{x})} = \frac{\sum_{\boldsymbol{x}_0} q_{t|0}(\tilde{\boldsymbol{x}} \mid \boldsymbol{x}_0) q_0(\boldsymbol{x}_0)}{\sum_{\boldsymbol{x}_0} q_{t|0}(\boldsymbol{x} \mid \boldsymbol{x}_0) q_0(\boldsymbol{x}_0)}. \label{eq:ratio-expand}
\end{equation}
按照 \cite[§2]{austin2021d3pm} 的 $x_0$ 参数化思路，引入神经网络 $p_{0|t}^\theta(\boldsymbol{x}_0 \mid \boldsymbol{x}_t)$
来近似不可知的 $q_0$，得 \cite[§3.1]{campbell2022ctmc}：
\begin{equation}
\hat{R}_t^\theta(\boldsymbol{x}, \tilde{\boldsymbol{x}}) = R_t(\tilde{\boldsymbol{x}}, \boldsymbol{x}) \sum_{\boldsymbol{x}_0} \frac{q_{t|0}(\tilde{\boldsymbol{x}} \mid \boldsymbol{x}_0)}{q_{t|0}(\boldsymbol{x} \mid \boldsymbol{x}_0)} p_{0|t}^\theta(\boldsymbol{x}_0 \mid \boldsymbol{x}), \quad \boldsymbol{x} \neq \tilde{\boldsymbol{x}}. \label{eq:reverse-approx}
\end{equation}
对角线元素由行和归零条件确定：
\begin{equation}
\hat{R}_t^\theta(\boldsymbol{x}, \boldsymbol{x}) = -\sum_{\tilde{\boldsymbol{x}} \neq \boldsymbol{x}} \hat{R}_t^\theta(\boldsymbol{x}, \tilde{\boldsymbol{x}}). \label{eq:reverse-diag}
\end{equation}
当 $p_{0|t}^\theta$ 精确逼近真实后验 $q_{0|t}$ 时，\eqref{eq:reverse-approx} 退化为 \eqref{eq:reverse-exact}。

\subsection{去噪网络的结构}

去噪网络 $p_{0|t}^\theta(\boldsymbol{x}_0 \mid \boldsymbol{x}_t)$ 接收带噪数据和时间 $t$，
输出 $K$ 维概率向量。各维度在给定 $\boldsymbol{x}_t$ 条件下条件独立：
\begin{equation}
p_{0|t}^\theta(\boldsymbol{x}_0^{1:D} \mid \boldsymbol{x}_t^{1:D}) = \prod_{d=1}^D p_{0|t}^{\theta,d}(x_0^d \mid \boldsymbol{x}_t^{1:D}). \label{eq:net-factorization}
\end{equation}
$\boldsymbol{x}_t$ 本身包含跨维度信息，网络需通过自注意力等机制捕获维度间的依赖关系。
交叉熵训练目标为
\begin{equation}
L_{\text{CE}}(\theta) = \mathbb{E}_{t, q_{t|0}(\boldsymbol{x}_t \mid \boldsymbol{x}_0) q_0(\boldsymbol{x}_0)}\left[-\log p_{0|t}^\theta(\boldsymbol{x}_0 \mid \boldsymbol{x}_t)\right]. \label{eq:ce-loss}
\end{equation}
完整训练目标详见 \ref{sec:elbo} 节。

\subsection{反向过程的维度分解}
\label{sec:reverse-dim-decomp}

反向过程继承前向的稀疏性结构 \cite[§4.2 Proposition 3]{campbell2022ctmc}：
\begin{equation}
\hat{R}_t^{1:D}(\boldsymbol{x}^{1:D}, \tilde{\boldsymbol{x}}^{1:D}) = \sum_{d=1}^D \hat{R}_t^d(\boldsymbol{x}^{1:D}, \tilde{x}^d) \;\delta_{\tilde{\boldsymbol{x}}^{1:D \setminus d}, \boldsymbol{x}^{1:D \setminus d}}. \label{eq:reverse-dim-decomp}
\end{equation}
其中偏反向速率为
\begin{equation}
\hat{R}_t^d(\boldsymbol{x}^{1:D}, \tilde{x}^d) = R_t^d(\tilde{x}^d, x^d) \sum_{x_0^d} \frac{q_{t|0}(\tilde{x}^d \mid x_0^d)}{q_{t|0}(x^d \mid x_0^d)} \; q_{0|t}(x_0^d \mid \boldsymbol{x}^{1:D}). \label{eq:reverse-dim-rate}
\end{equation}
$q_{0|t}(x_0^d \mid \boldsymbol{x}^{1:D})$ 是\textbf{跨维度耦合}的，这是生成模型需要学习的核心。
参数化形式为
\begin{equation}
\hat{R}_t^{\theta,d}(\boldsymbol{x}^{1:D}, \tilde{x}^d) = R_t^d(\tilde{x}^d, x^d) \sum_{x_0^d} \frac{q_{t|0}(\tilde{x}^d \mid x_0^d)}{q_{t|0}(x^d \mid x_0^d)} \; p_{0|t}^{\theta,d}(x_0^d \mid \boldsymbol{x}^{1:D}). \label{eq:reverse-dim-param}
\end{equation}
单次网络前向即可得到整个 $D \times (K-1)$ 维的反向速率张量。

\subsection{与离散时间反向过程的联系}

D3PM 的 $x_0$ 参数化 \cite[§2]{austin2021d3pm} 为
\begin{equation}
p_\theta(\boldsymbol{x}_{t-1} \mid \boldsymbol{x}_t) \propto \sum_{\boldsymbol{x}_0} q(\boldsymbol{x}_{t-1}, \boldsymbol{x}_t \mid \boldsymbol{x}_0) \, \tilde{p}_\theta(\boldsymbol{x}_0 \mid \boldsymbol{x}_t). \label{eq:d3pm-reverse}
\end{equation}
连续时间参数化 \eqref{eq:reverse-approx} 可视为 \eqref{eq:d3pm-reverse} 在 $\Delta t \to 0$ 时的极限。
两者的核心思想一致，为 \ref{sec:elbo} 节的极限推导奠定了基础。

% ============================================================
\section{连续时间 ELBO}
\label{sec:elbo}

训练目标是最小化生成分布与数据分布之间的差距，通过最大化证据下界实现。

\subsection{变分下界}

变分方法引入前向过程 $q$ 作为辅助分布：
\begin{equation}
-\log p_0^\theta(x_0) \leq \mathbb{E}_{q(x_{1:K} \mid x_0)}\left[-\log \frac{p_{0:K}^\theta(x_{0:K})}{q(x_{1:K} \mid x_0)}\right] =: \mathcal{L}_{\text{DT}}(\theta). \label{eq:elbo-basic}
\end{equation}
离散时间 ELBO 分解为三项 \cite[§2]{austin2021d3pm}：
\begin{equation}
\begin{aligned}
\mathcal{L}_{\text{DT}}(\theta) = \mathbb{E}_{p_{\text{data}}(x_0)}\Big[
&\underbrace{D_{\text{KL}}(q(x_T \mid x_0) \parallel p_{\text{ref}}(x_T))}_{L_T} \\
&+ \sum_{t=2}^T \underbrace{\mathbb{E}_{q(x_t \mid x_0)}\big[D_{\text{KL}}(q(x_{t-1} \mid x_t, x_0) \parallel p_\theta(x_{t-1} \mid x_t))\big]}_{L_{t-1}} \\
&+ \underbrace{\big(-\mathbb{E}_{q(x_1 \mid x_0)}[\log p_\theta(x_0 \mid x_1)]\big)}_{L_0}
\Big]. \label{eq:elbo-discrete}
\end{aligned}
\end{equation}

\subsection{用速率矩阵表示离散 ELBO 的局部项}

将时间区间 $[0, T]$ 均匀划分为 $K$ 步，步长 $\Delta t = T/K$。
由 \cite[Appendix B.2]{campbell2022ctmc}，转移核的无穷小展开为
\begin{align}
p_{k|k+1}^\theta(x_k \mid x_{k+1}) &= \delta_{x_k, x_{k+1}} + \hat{R}_k^\theta(x_{k+1}, x_k) \Delta t + o(\Delta t), \label{eq:p-expand} \\
q_{k+1|k}(x_{k+1} \mid x_k) &= \delta_{x_k, x_{k+1}} + R_k(x_k, x_{k+1}) \Delta t + o(\Delta t). \label{eq:q-expand}
\end{align}
取对数并展开：
\begin{equation}
\log p_{k|k+1}^\theta(x_k \mid x_{k+1}) = \delta_{x_k, x_{k+1}} \hat{R}_k^\theta(x_k, x_k) \Delta t + (1 - \delta_{x_k, x_{k+1}}) \log\big(\hat{R}_k^\theta(x_{k+1}, x_k) \Delta t + o(\Delta t)\big) + o(\Delta t). \label{eq:log-p-expand}
\end{equation}
将 $\log(\hat{R}_k^\theta \Delta t + o(\Delta t)) = \log \Delta t + \log \hat{R}_k^\theta + o(1)$ 代入，
$\log \Delta t$ 项与 $\theta$ 无关，经期望操作后仅贡献一个加性常数。
整理得：
\begin{equation}
L_k = -\mathbb{E}_{q_k(x_k)}\Big[ \hat{R}_k^\theta(x_k, x_k) \Delta t + \sum_{x_{k+1} \neq x_k} R_k(x_k, x_{k+1}) \Delta t \log \hat{R}_k^\theta(x_{k+1}, x_k) + o(\Delta t) \Big] + C. \label{eq:l-k}
\end{equation}

\subsection{连续极限下的 ELBO}

取 $\Delta t \to 0$（$K \to \infty$）的极限，离散求和转化为积分，
得到连续时间 ELBO 的闭式表达 \cite[§3.2 Proposition 2]{campbell2022ctmc}：
\begin{equation}
\mathcal{L}_{\text{CT}}(\theta) = T \,\mathbb{E}_{t \sim \mathcal{U}(0,T),\; q_t(x),\; r_t(\tilde{x} \mid x)}\Big[ \Big\{\sum_{x' \neq x} \hat{R}_t^\theta(x, x')\Big\} - \mathcal{Z}^t(x) \log \big(\hat{R}_t^\theta(\tilde{x}, x)\big) \Big] + C. \label{eq:elbo-ct}
\end{equation}
其中
\begin{equation}
\mathcal{Z}^t(x) = \sum_{x' \neq x} R_t(x, x'), \qquad
r_t(\tilde{x} \mid x) = \frac{R_t(x, \tilde{x})}{\mathcal{Z}^t(x)}\; (\tilde{x} \neq x). \label{eq:z-and-r}
\end{equation}
$\mathcal{Z}^t(x)$ 是总跳出速率，$r_t(\tilde{x} \mid x)$ 是跳转目的地的条件分布。
式 \eqref{eq:elbo-ct} 的两项形成一种平衡：反向过程不应无理由跳转（第一项的惩罚），
但当正向过程发生跳转时须学会从目的地跳回（第二项的鼓励）。

\subsection{高效单次前向近似}

式 \eqref{eq:elbo-ct} 理论上需要两次网络前向。\cite[Appendix C.4]{campbell2022ctmc} 提出近似：
\begin{equation}
\mathcal{L}_{\text{eCT}}(\theta) = T \,\mathbb{E}_{t \sim \mathcal{U}(0,T),\; q_t(x),\; r_t(\tilde{x} \mid x)}\Big[ \Big\{\sum_{x' \neq \tilde{x}} \hat{R}_t^\theta(\tilde{x}, x')\Big\} - \mathcal{Z}^t(x) \log \big(\hat{R}_t^\theta(\tilde{x}, x)\big) \Big] + C. \label{eq:elbo-ect}
\end{equation}
两项共享同一个 $\tilde{x}$，计算量减半且方差降低。

\subsection{辅助去噪损失}

为稳定训练，\cite[Appendix D]{campbell2022ctmc} 建议叠加直接去噪损失：
\begin{equation}
L_{\text{ll}}(\theta) = T \,\mathbb{E}_{t \sim \mathcal{U}(0,T),\; p_{\text{data}}(x_0),\; q_{t|0}(x \mid x_0)}\big[ -\log p_{0|t}^\theta(x_0 \mid x) \big]. \label{eq:ll-loss}
\end{equation}
当 $p_{0|t}^\theta = q_{0|t}$ 时 $L_{\text{ll}}$ 达到最小值 \cite[Proposition 8]{campbell2022ctmc}。
完整目标为二者加权组合：
\begin{equation}
\min_\theta \mathcal{L}_{\text{CT}}(\theta) + \lambda L_{\text{ll}}(\theta). \label{eq:combined-loss}
\end{equation}

\subsection{高维数据的因子化 ELBO}

利用维度分解 \cite[Appendix C.3 Proposition 7]{campbell2022ctmc}：
\begin{equation}
\begin{aligned}
\mathcal{L}_{\text{CT}} = T \,\mathbb{E}_{t, p_{\text{data}}(\boldsymbol{x}_0^{1:D}), q_{t|0}(\boldsymbol{x}^{1:D} \mid \boldsymbol{x}_0^{1:D})}
&\Big[ \sum_{d=1}^D \sum_{x'^d \neq x^d} \hat{R}_t^{\theta d}(\boldsymbol{x}^{1:D}, x'^d) \Big] \\
&- T \,\mathbb{E}_{t, p_{\text{data}}(\boldsymbol{x}_0^{1:D}), \psi_t(\tilde{\boldsymbol{x}}^{1:D} \mid \boldsymbol{x}_0^{1:D})}
\Big[ \sum_{d=1}^D \sum_{x^d \neq \tilde{x}^d} \phi_t(x^d \mid \tilde{\boldsymbol{x}}^{1:D}, \boldsymbol{x}_0^{1:D}) \mathcal{Z}^t(\tilde{\boldsymbol{x}}^{1:D/d} \circ x^d) \log \hat{R}_t^{\theta d}(\tilde{\boldsymbol{x}}^{1:D}, x^d) \Big] + C. \label{eq:elbo-factorized}
\end{aligned}
\end{equation}
将计算复杂度降到 $O(DK)$，避免了全空间 $\mathcal{S}^D$ 上的求和。

% ============================================================
\section{反向采样算法}
\label{sec:sampling}

给定反向速率矩阵 $\hat{R}_t^\theta$，需要从中高效生成样本。
Gillespie 精确模拟每步只改变一个维度，在高维数据上不可行。

\subsection{Tau-Leaping 算法}

Tau-leaping 假设速率矩阵在 $[t-\tau, t]$ 内恒定，批量处理所有可能跳转 \cite[§4.3]{campbell2022ctmc}：
\begin{equation}
P_{ds} \sim \text{Poisson}\Big( \tau \,\hat{R}_t^{\theta,d}(\boldsymbol{x}_t^{1:D}, s) \Big). \label{eq:poisson}
\end{equation}
更新规则为
\begin{equation}
\boldsymbol{x}_{t-\tau}^{1:D} = \boldsymbol{x}_t^{1:D} + \sum_{d=1}^D \sum_{s \neq x_t^d} P_{ds} \, (s - x_t^d) \, \boldsymbol{e}^d. \label{eq:tau-update}
\end{equation}
对于类别数据，同一维度的多次跳转无意义，当 $\sum_{s} P_{ds} > 1$ 时拒绝该维度更新：
\begin{equation}
\boldsymbol{x}_{t-\tau}^d \leftarrow \boldsymbol{x}_t^d. \label{eq:tau-reject}
\end{equation}
总网络前向次数为 $\text{NFE} = T / \tau$ \cite[§4.3]{campbell2022ctmc}。
步长 $\tau$ 控制效率与精度的权衡。

\subsection{Predictor-Corrector 采样器}

\cite[§4.4 Proposition 4]{campbell2022ctmc} 发现 $R_t^c = R_t + \hat{R}_t$ 以 $q_t$ 为平稳分布：
\begin{equation}
R_t^c = R_t + \hat{R}_t. \label{eq:corrector-rate}
\end{equation}
实际流程 \cite[Appendix F.2]{campbell2022ctmc} 交替执行预测步（$\hat{R}_t^\theta$ tau-leaping）
和校正步（$R_t^{c\theta} = \hat{R}_t^\theta + R_t$ tau-leaping）。
校正步将近似误差向 $q_t$ "拉回"。CIFAR-10 上的实验结果见表 \ref{tab:cifar10}：

\begin{table}[htbp]
\centering
\caption{CIFAR-10 上 Predictor-Corrector 采样器的性能对比 \cite{campbell2022ctmc}}
\label{tab:cifar10}
\begin{tabular}{lcc}
\toprule
方法 & IS ($\uparrow$) & FID ($\downarrow$) \\
\midrule
D3PM Gaussian & 8.56 & 7.34 \\
$\tau$LDR-0（无校正步） & 8.74 & 8.10 \\
$\tau$LDR-10（10 校正步） & 9.49 & 3.74 \\
DDPM（连续状态空间） & 9.46 & 3.17 \\
\bottomrule
\end{tabular}
\end{table}

\subsection{误差界}

\cite[§4.5 Theorem 1]{campbell2022ctmc} 给出生成分布与数据分布之间的总变差上界：
\begin{equation}
\begin{aligned}
\|\mathcal{L}(y_0) - p_{\text{data}}\|_{\text{TV}} \le\; &3MT \\
&+ \Big\{ \big(|R| S D C_1\big)^2 + \frac{1}{2} C_2 (M + C_1 S D |R|) \Big\} \tau T \\
&+ 2 \exp\!\Big\{ -\frac{T \log^2 2}{t_{\text{mix}} \log 4D} \Big\}. \label{eq:error-bound}
\end{aligned}
\end{equation}
三项依次对应反向速率近似误差、tau-leaping 离散化误差和初始分布不匹配误差。
最重要的结论是总误差关于维度 $D$ 的增长是 $\mathcal{O}(D^2)$ 而非 $\mathcal{O}(S^D)$，
这一多项式可扩展性是 CTMC 框架的关键优势。

% ============================================================
\section{连续时间与离散时间模型的统一}
\label{sec:ud3}

前几节分别从离散时间和连续时间两个角度构建了扩散模型。
\cite{zhao2025ud3} 提出的统一框架（UD3）使得两类模型共享完全相同的前向过程、反向参数化和采样流程。

\subsection{幂等速率矩阵与统一的数学根源}

统一的关键在于选择 \cite[§2.1]{zhao2025ud3}：
\begin{equation}
R_b = \mathbf{1} \mathbf{m}^\top - I, \label{eq:idempotent}
\end{equation}
其中 $\mathbf{m}$ 是平稳分布。$(-R_b)^2 = -R_b$ 的幂等性使矩阵指数坍缩为两项：
\begin{equation}
\exp(\alpha R_b) = e^{-\alpha} I + (1 - e^{-\alpha}) \mathbf{1} \mathbf{m}^\top. \label{eq:exp-collapse}
\end{equation}
令 $\overline{\alpha}_{t|s} = e^{-\overline{\beta}_{t|s}}$，则连续时间转移概率
\begin{equation}
\overline{Q}_{t|s} = \overline{\alpha}_{t|s} I + (1 - \overline{\alpha}_{t|s}) \mathbf{1} \mathbf{m}^\top \label{eq:q-bar}
\end{equation}
与离散时间形式完全相同，开启了统一的大门。

\subsection{闭式后验}

由 \cite[§2.2 Proposition 2.1]{zhao2025ud3}：
\begin{equation}
q(\boldsymbol{x}_s \mid \boldsymbol{x}_t, \boldsymbol{x}_0) =
\begin{cases}
\text{Cat}\big(\boldsymbol{x}_s; (1 - \lambda_{t|s})\boldsymbol{x}_t + \lambda_{t|s} \mathbf{m}\big), & \boldsymbol{x}_t = \boldsymbol{x}_0, \\[6pt]
\text{Cat}\big(\boldsymbol{x}_s; (1 - \mu_{t|s})\boldsymbol{x}_0 + \mu_{t|s}\overline{\alpha}_{t|s}\boldsymbol{x}_t + \mu_{t|s}(1 - \overline{\alpha}_{t|s})\mathbf{m}\big), & \boldsymbol{x}_t \neq \boldsymbol{x}_0.
\end{cases} \label{eq:posterior}
\end{equation}
系数为
\begin{equation}
\lambda_{t|s} = \frac{(1 - \overline{\alpha}_s)(1 - \overline{\alpha}_{t|s}) \langle \mathbf{m}, \boldsymbol{x}_t \rangle}{\overline{\alpha}_t + (1 - \overline{\alpha}_t) \langle \mathbf{m}, \boldsymbol{x}_t \rangle},
\qquad
\mu_{t|s} = \frac{1 - \overline{\alpha}_s}{1 - \overline{\alpha}_t}. \label{eq:lambda-mu}
\end{equation}
后验从 $\mathcal{O}(K)$ 求和简化为两个分情况讨论的闭式。

\subsection{闭式反向分布}

按照 $x_0$ 参数化，反向分布通过对 $\boldsymbol{x}_0$ 边缘化：
\begin{equation}
p_\theta(\boldsymbol{x}_s \mid \boldsymbol{x}_t) = \sum_{\boldsymbol{x}_0} q(\boldsymbol{x}_s \mid \boldsymbol{x}_t, \boldsymbol{x}_0) \, p_{0|t}^\theta(\boldsymbol{x}_0 \mid \boldsymbol{x}_t). \label{eq:p-reverse-sum}
\end{equation}
\cite[§2.2 Proposition 2.2]{zhao2025ud3} 证明该求和化简为三项混合分布：
\begin{equation}
p_\theta(\boldsymbol{x}_s \mid \boldsymbol{x}_t) = \text{Cat}\Big(\boldsymbol{x}_s; \;
(1 - \mu_{t|s}) \, f_t^\theta(\boldsymbol{x}_t) + (\mu_{t|s}\overline{\alpha}_{t|s} + \gamma_{t|s}^\theta) \, \boldsymbol{x}_t + (\mu_{t|s}(1 - \overline{\alpha}_{t|s}) - \gamma_{t|s}^\theta) \, \mathbf{m} \Big), \label{eq:p-reverse-closed}
\end{equation}
其中 $f_t^\theta(\boldsymbol{x}_t) = p_{0|t}^\theta(\boldsymbol{x}_0 \mid \boldsymbol{x}_t)$，
修正系数
\begin{equation}
\gamma_{t|s}^\theta := (\mu_{t|s} - \lambda_{t|s} - \mu_{t|s}\overline{\alpha}_{t|s}) \, \langle f_t^\theta(\boldsymbol{x}_t), \boldsymbol{x}_t \rangle. \label{eq:gamma}
\end{equation}
重参数化采样形式为
\begin{equation}
\boldsymbol{x}_{s|t} = \delta_{1, b_{s|t}} \, \tilde{\boldsymbol{x}}_{0|t} + \delta_{2, b_{s|t}} \, \boldsymbol{x}_t + \delta_{3, b_{s|t}} \, \mathbf{m},
\qquad
b_{s|t} \sim \text{Cat}\big([1 - \mu_{t|s}, \; \mu_{t|s}\overline{\alpha}_{t|s} + \gamma_{t|s}^\theta, \; \mu_{t|s}(1 - \overline{\alpha}_{t|s}) - \gamma_{t|s}^\theta]\big). \label{eq:reparam-sample}
\end{equation}
对任意时间调度均可使用，无需 tau-leaping 近似。

\subsection{简化连续时间变分下界}

UD3 引入比率估计器 $g_t^\theta$ 来近似边缘概率比 \cite[§3.1]{zhao2025ud3}：
\begin{equation}
g_t^\theta(x \mid y) := \sum_{x_0} \frac{q_{t|0}(x \mid x_0)}{q_{t|0}(y \mid x_0)} p_{0|t}^\theta(x_0 \mid y) \approx \frac{q_t(x)}{q_t(y)}. \label{eq:g-def}
\end{equation}
利用闭式结构向量化：
\begin{equation}
g_t^\theta(\cdot \mid \boldsymbol{y}) = \left[ \Big(1 - \frac{\overline{\alpha}_{t|0} \langle f_t^\theta(\boldsymbol{y}), \boldsymbol{y} \rangle}{\overline{\alpha}_{t|0} + (1 - \overline{\alpha}_{t|0})\langle \boldsymbol{y}, \mathbf{m} \rangle}\Big) \mathbf{m} + \frac{\overline{\alpha}_{t|0}}{1 - \overline{\alpha}_{t|0}} f_t^\theta(\boldsymbol{y}) \right] \odot \frac{\mathbf{1} - \boldsymbol{y}}{\langle \boldsymbol{y}, \mathbf{m} \rangle} + \boldsymbol{y}. \label{eq:g-closed}
\end{equation}
由此得到单次前向版本的连续时间 ELBO \cite[§3.1 Proposition 3.3]{zhao2025ud3}：
\begin{equation}
\mathcal{L}_{\text{CT}}^{\text{UD3}}(\theta) = T \,\mathbb{E}_{t, \boldsymbol{x}^{1:D} \sim q_{t|0}}\Big[
\sum_{d=1}^D r_t^d(x^d \mid \cdot)^\top \Big(
g_t^{\theta,d}(\cdot \mid \boldsymbol{x}^{1:D}) - \frac{q_{t|0}(\cdot \mid x_0^d) \odot \log g_t^{\theta,d}(\cdot \mid \boldsymbol{x}^{1:D})}{q_{t|0}(x^d \mid x_0^d)}
\Big)
\Big]. \label{eq:elbo-ud3}
\end{equation}
训练计算成本减半且方差更低。

\subsection{统一的反向采样与双向收益}

UD3 框架下离散时间和连续时间共享完全相同的反向采样流程。
对连续时间模型，使用 \eqref{eq:p-reverse-closed} 的闭式反向分布可进行精确采样：
\begin{equation}
\boldsymbol{x}_{t_{i-1} \mid t_i} \sim p_\theta(\boldsymbol{x}_{t_{i-1}} \mid \boldsymbol{x}_{t_i}), \quad i = n, n-1, \dots, 1. \label{eq:unified-sample}
\end{equation}
这从根本上解决了 tau-leaping 的离散化误差。对离散时间模型，
可利用连续时间框架的 MCMC 校正子来改善样本质量 \cite[Appendix A.19]{zhao2025ud3}。
离散时间和连续时间并非两种不同模型，而是同一模型在不同时间离散化粒度下的实例。

% ---------- 参考文献 ----------
\bibliographystyle{plainnat}
\bibliography{references}

\end{document}
